% Source slide 83
\section[4.6 Rank]{4.6 Rank}
\chaptercard{4}{Vector Spaces}{4.6}{RANK}

% Source slide 84
\begin{frame}[shrink=8]{THE ROW SPACE}
\begin{itemize}
\item If $A$ is an $m\times n$ matrix, each row of $A$ has $n$ entries and thus can be identified with a vector in $\R^n$.
\item The set of all linear combinations of the row vectors is called the \textbf{row space} of $A$ and is denoted by $\Row A$.
\item Each row has $n$ entries, so $\Row A$ is a subspace of $\R^n$.
\item Since the rows of $A$ are identified with the columns of $A^T$, we could also write $\Col A^T$ in place of $\Row A$.
\end{itemize}
\end{frame}

% Source slide 85
\begin{frame}[shrink=8]{THE ROW SPACE}
\begin{itemize}
\item \sourcehead{Theorem 13:} If two matrices $A$ and $B$ are row equivalent, then their row spaces are the same. If $B$ is in echelon form, the nonzero rows of $B$ form a basis for the row space of $A$ as well as for that of $B$.
\item \textbf{Proof:} If $B$ is obtained from $A$ by row operations, the rows of $B$ are linear combinations of the rows of $A$.
\item It follows that any linear combination of the rows of $B$ is automatically a linear combination of the rows of $A$.
\end{itemize}
\end{frame}

% Source slide 86
\begin{frame}{THE ROW SPACE}
\begin{itemize}
\item Thus the row space of $B$ is contained in the row space of $A$.
\item Since row operations are reversible, the same argument shows that the row space of $A$ is a subset of the row space of $B$.
\item So the two row spaces are the same.
\end{itemize}
\end{frame}

% Source slide 87
\begin{frame}[shrink=7]{THE ROW SPACE}
\begin{itemize}
\item If $B$ is in echelon form, its nonzero rows are linearly independent because no nonzero row is a linear combination of the nonzero rows below it. (Apply Theorem 4 to the nonzero rows of $B$ in reverse order, with the first row last).
\item Thus the nonzero rows of $B$ form a basis of the (common) row space of $B$ and $A$.
\end{itemize}
\end{frame}

% Source slide 88
\begin{frame}[shrink=9]{THE ROW SPACE}
\begin{itemize}
\item \textbf{Example 2:} Find bases for the row space, the column space, and the null space of the matrix
\[
A=\begin{bmatrix}
-2&-5&8&0&-17\\
1&3&-5&1&5\\
3&11&-19&7&1\\
1&7&-13&5&-3
\end{bmatrix}.
\]
\item \textbf{Solution:} To find bases for the row space and the column space, row reduce $A$ to an echelon form:
\end{itemize}
\end{frame}

% Source slide 89
\begin{frame}[shrink=10]{THE ROW SPACE}
\[
A\sim B=\begin{bmatrix}
1&3&-5&1&5\\
0&1&-2&2&-7\\
0&0&0&-4&20\\
0&0&0&0&0
\end{bmatrix}.
\]
\begin{itemize}
\item By Theorem 13, the first three rows of $B$ form a basis for the row space of $A$ (as well as for the row space of $B$).
\item Thus
\[
\text{Basis for }\Row A:\{(1,3,-5,1,5),(0,1,-2,2,-7),(0,0,0,-4,20)\}.
\]
\end{itemize}
\end{frame}

% Source slide 90
\begin{frame}[shrink=9]{THE ROW SPACE}
\begin{itemize}
\item For the column space, observe from $B$ that the pivots are in columns 1, 2, and 4.
\item Hence columns 1, 2, and 4 of $A$ (not $B$) form a basis for $\Col A$:
\[
\text{Basis for }\Col A:
\left\{
\begin{bmatrix}-2\\1\\3\\1\end{bmatrix},
\begin{bmatrix}-5\\3\\11\\7\end{bmatrix},
\begin{bmatrix}0\\1\\7\\5\end{bmatrix}
\right\}.
\]
\item Notice that any echelon form of $A$ provides (in its nonzero rows) a basis for $\Row A$ and also identifies the pivot columns of $A$ for $\Col A$.
\end{itemize}
\end{frame}

% Source slide 91
\begin{frame}{THE ROW SPACE}
\begin{itemize}
\item However, for $\Nul A$, we need the \emph{reduced echelon form}.
\item Further row operations on $B$ yield
\[
A\sim B\sim C=\begin{bmatrix}
1&0&1&0&1\\
0&1&-2&0&3\\
0&0&0&1&-5\\
0&0&0&0&0
\end{bmatrix}.
\]
\end{itemize}
\end{frame}

% Source slide 92
\begin{frame}[shrink=8]{THE ROW SPACE}
\begin{itemize}
\item The equation $A\vect x=0$ is equivalent to $C\vect x=0$, that is,
\begin{align*}
x_1+x_3+x_5&=0,\\
x_2-2x_3+3x_5&=0,\\
x_4-5x_5&=0.
\end{align*}
\item So $x_1=-x_3-x_5$, $x_2=2x_3-3x_5$, $x_4=5x_5$, with $x_3$ and $x_5$ free variables.
\end{itemize}
\end{frame}

% Source slide 93
\begin{frame}[shrink=8]{THE ROW SPACE}
\begin{itemize}
\item The calculations show that
\[
\text{Basis for }\Nul A:
\left\{
\begin{bmatrix}-1\\2\\1\\0\\0\end{bmatrix},
\begin{bmatrix}-1\\-3\\0\\5\\1\end{bmatrix}
\right\}.
\]
\item Observe that, unlike the basis for $\Col A$, the bases for $\Row A$ and $\Nul A$ have no simple connection with the entries in $A$ itself.
\end{itemize}
\end{frame}

% Source slide 94
\begin{frame}[shrink=11]{THE RANK THEOREM}
\begin{itemize}
\item \textbf{Definition:} The \textbf{rank} of $A$ is the dimension of the column space of $A$.
\item Since $\Row A$ is the same as $\Col A^T$, the dimension of the row space of $A$ is the rank of $A^T$.
\item The dimension of the null space is sometimes called the \textbf{nullity} of $A$.
\item \sourcehead{Theorem 14:} The dimensions of the column space and the row space of an $m\times n$ matrix $A$ are equal. This common dimension, the rank of $A$, also equals the number of pivot positions in $A$ and satisfies the equation
\[
\operatorname{rank}A+\dim\Nul A=n.
\]
\end{itemize}
\end{frame}

% Source slide 95
\begin{frame}[shrink=9]{THE RANK THEOREM}
\begin{itemize}
\item \textbf{Proof:} By Theorem 6, $\operatorname{rank}A$ is the number of pivot columns in $A$.
\item Equivalently, $\operatorname{rank}A$ is the number of pivot positions in an echelon form $B$ of $A$.
\item Since $B$ has a nonzero row for each pivot, and since these rows form a basis for the row space of $A$, the rank of $A$ is also the dimension of the row space.
\item The dimension of $\Nul A$ equals the number of free variables in the equation $A\vect x=0$.
\item Expressed another way, the dimension of $\Nul A$ is the number of columns of $A$ that are \emph{not} pivot columns.
\end{itemize}
\end{frame}

% Source slide 96
\begin{frame}[shrink=8]{THE RANK THEOREM}
\begin{itemize}
\item (It is the number of these columns, not the columns themselves, that is related to $\Nul A$).
\item Obviously,
\[
\underbrace{\begin{gathered}\text{number of}\\\text{pivot columns}\end{gathered}}_{\vphantom{X}}
+
\underbrace{\begin{gathered}\text{number of}\\\text{nonpivot columns}\end{gathered}}_{\vphantom{X}}
=
\underbrace{\begin{gathered}\text{number of}\\\text{columns}\end{gathered}}_{\vphantom{X}}.
\]
\item This proves the theorem.
\end{itemize}
\end{frame}

% Source slide 97
\begin{frame}[shrink=8]{THE RANK THEOREM}
\begin{itemize}
\item \textbf{Example 3:}
\begin{itemize}
\item[a.] If $A$ is a $7\times9$ matrix with a two-dimensional null space, what is the rank of $A$?
\item[b.] Could a $6\times9$ matrix have a two-dimensional null space?
\end{itemize}
\item \textbf{Solution:}
\begin{itemize}
\item[a.] Since $A$ has 9 columns, $(\operatorname{rank}A)+2=9$, and hence $\operatorname{rank}A=7$.
\item[b.] No. If a $6\times9$ matrix, call it $B$, has a two-dimensional null space, it would have to have rank 7, by the Rank Theorem.
\end{itemize}
\end{itemize}
\end{frame}

% Source slide 98
\begin{frame}[shrink=9]{THE INVERTIBLE MATRIX THEOREM}
\textbf{(CONTINUED)}
\begin{itemize}
\item But the columns of $B$ are vectors in $\R^6$, and so the dimension of $\Col B$ cannot exceed 6; that is, $\operatorname{rank}B$ cannot exceed 6.
\item \textbf{Theorem:} Let $A$ be an $n\times n$ matrix. Then the following statements are each equivalent to the statement that $A$ is an invertible matrix.
\begin{itemize}
\item[m.] The columns of $A$ form a basis of $\R^n$.
\item[n.] $\Col A=\R^n$.
\item[o.] $\dim\Col A=n$.
\item[p.] $\operatorname{rank}A=n$.
\end{itemize}
\end{itemize}
\end{frame}

% Source slide 99
\begin{frame}[shrink=8]{RANK AND THE INVERTIBLE MATRIX THEOREM}
\begin{itemize}
\item[q.] $\Nul A=\{0\}$.
\item[r.] $\dim\Nul A=0$.
\item \textbf{Proof:} Statement (m) is logically equivalent to statements (e) and (h) regarding linear independence and spanning.
\item The other five statements are linked to the earlier ones of the theorem by the following chain of almost trivial implications:
\[
(g)\Rightarrow(n)\Rightarrow(o)\Rightarrow(p)\Rightarrow(r)\Rightarrow(q)\Rightarrow(d).
\]
\end{itemize}
\end{frame}

% Source slide 100
\begin{frame}[shrink=9]{RANK AND THE INVERTIBLE MATRIX THEOREM}
\begin{itemize}
\item Statement (g), which says that the equation $A\vect x=\vect b$ has at least one solution for each $\vect b$ in $\R^n$, implies (n), because $\Col A$ is precisely the set of all $\vect b$ such that the equation $A\vect x=\vect b$ is consistent.
\item The implications $(n)\Rightarrow(o)\Rightarrow(p)$ follow from the definitions of dimension and rank.
\item If the rank of $A$ is $n$, the number of columns of $A$, then $\dim\Nul A=0$, by the Rank Theorem, and so $\Nul A=\{0\}$.
\end{itemize}
\end{frame}

% Source slide 101
\begin{frame}{RANK AND THE INVERTIBLE MATRIX THEOREM}
\begin{itemize}
\item Thus $(p)\Rightarrow(r)\Rightarrow(q)$.
\item Also, (q) implies that the equation $A\vect x=0$ has only the trivial solution, which is statement (d).
\item Since statements (d) and (g) are already known to be equivalent to the statement that $A$ is invertible, the proof is complete.
\end{itemize}
\end{frame}
